\section{前期研究：多算为何收效甚微}

本研究始于作者在东海赤潮巡测场景下对机载算力分配的初步探索：当航行器携带3天前的卫星叶绿素浓度图下水，在“继续计算”“边走边测”与“停下复测”之间分配有限作业时间时，多算能否带来更好的观测收益\cite{lou2014diurnal,mnr2024bulletin}？我们用多年卫星叶绿素数据标定旧图误差并构造48个想象海洋（32个用于搜索、16个留出不参与搜索）\cite{beckers2003dineof,gaspari1999correlation}，采用交叉熵方法搜索巡测航线\cite{deboer2005cem}，以航行器报告的最浓位置与真实最浓处之间对数叶绿素浓度差的时间平均（遗憾）为指标，考察了推演次数与用时策略的作用。

前期研究的第一个发现是，在陈旧背景图上大幅增加推演计算，只会使规划器在自身搜索过的想象海洋上给出的自评持续虚高，而航线在真实海洋中的价值早已饱和\cite{smith2006curse,vanhasselt2010double}。当每次决策的推演次数从64次增至2000次与20万次时，规划器自评从0.762升至0.833和0.917；但所选航线在真实海洋中的价值仅从0.749升至0.793和0.803，在约2000次推演后便饱和，且在不参与搜索的留出海洋上的评判始终紧贴真实价值（图\ref{fig:curse}）。这表明过多算力并未帮航行器看清真实赤潮位置，而是被用来迎合有限想象样本中的偶然高值噪声。

\begin{figure}[htbp]
\centering
\includegraphics[width=\linewidth]{fig_curse}
\caption{前期东海赤潮巡测研究揭示的机载过度计算现象。(a) 所选航线价值（对数叶绿素浓度）随每次决策想象推演次数（64至20万次）的变化：规划器搜索时的自评从0.762持续升至0.917，而真实价值从0.749升至0.793后在约2000次推演处饱和（20万次时为0.803），留出海洋上的独立评判紧贴真实价值；(b) 估计价值减真实价值的偏差：搜索自评的高估偏差随计算量单调膨胀，留出评判保持无偏，不确定性幅度设错仅引起整体平移。多算只让规划器的答案在旧信息上越来越“精致”，却不能提供区分真实环境的新信息。}
\label{fig:curse}
\end{figure}

前期研究的第二个发现是，在连续平滑的标量场巡测中，计算与观测的时间分配所能改变的结局空间十分有限。出发前预演选出的用时策略平均遗憾为0.265，而事后为每个任务挑最好的用时策略也只能降到0.243，用时分配本身的空间仅约0.022。对遗憾来源的拆解进一步显示，下水位置决定了约0.13的差距，信息新鲜度决定了约0.03，而用时分配仅占约0.02；尝试用前一天光流外推3天反而比直接用旧图误差增大0.018，说明逐日卫星图里没有可外推的整体移动。

这一前期探索促使作者把问题收紧：既然当多个可能世界都与旧图相容、却要求采取不同行动时，继续计算无法将它们区分开来，那么在强海流绕障这种选错一侧就会付出数小时代价的二元拓扑决策中，能否直接检验“是否存在另一个同样合理、却使当前选择明显受损的流场”，并以此决定何时该算、何时该测、何时该保守绕行？这正是本文接下来要解决的问题。
